\documentclass[twoside,11pt]{article}

% Any additional packages needed should be included after jmlr2e.
% Note that jmlr2e.sty includes epsfig, amssymb, natbib and graphicx,
% and defines many common macros, such as 'proof' and 'example'.
%
% It also sets the bibliographystyle to plainnat; for more information on
% natbib citation styles, see the natbib documentation, a copy of which
% is archived at http://www.jmlr.org/format/natbib.pdf

\usepackage{jmlr2e}
\usepackage{lipsum}
\usepackage{amsmath}
\usepackage{hyperref}
\usepackage{algorithm,algorithmic}

% Definitions of handy macros can go here

\newcommand{\dataset}{{\cal D}}
\newcommand{\fracpartial}[2]{\frac{\partial #1}{\partial  #2}}

\renewcommand{\algorithmicrequire}{\textbf{Input:}}
\renewcommand{\algorithmicensure}{\textbf{Output:}}

% Heading arguments are {volume}{year}{pages}{submitted}{published}{author-full-names}

\jmlrheading{1}{2023}{1-10}{1/09}{12/12}{Insert your full name}

% Short headings should be running head and authors last names

\ShortHeadings{Insert short title of a paper}{Insert your last name}
\firstpageno{1}

\begin{document}

\title{Paper title}

\author{\name Your full name \email your.email@stud.mif.vu.lt \\
\name Your full name \email your.email@stud.mif.vu.lt \\
\name Your full name \email your.email@stud.mif.vu.lt \\
\name Your full name \email your.email@stud.mif.vu.lt \\
       \addr Data Science study programme\\
       Faculty of Mathematics and Informatics}

\editor{Jurgita Markevi\v{c}i\={u}t\.{e}}

\maketitle

\begin{abstract}%   <- trailing '%' for backward compatibility of .sty file
This paper describes the mixtures-of-trees model, a probabilistic
model for discrete multidimensional domains.  Mixtures-of-trees
generalize the probabilistic trees of \citet{chow:68}
in a different and complementary direction to that of Bayesian networks.
We present efficient algorithms for learning mixtures-of-trees
models in maximum likelihood and Bayesian frameworks.
We also discuss additional efficiencies that can be
obtained when data are ``sparse,'' and we present data
structures and algorithms that exploit such sparseness.
Experimental results demonstrate the performance of the
model for both density estimation and classification.
We also discuss the sense in which tree-based classifiers
perform an implicit form of feature selection, and demonstrate
a resulting insensitivity to irrelevant attributes.
\end{abstract}

\begin{keywords}
  Bayesian Networks, Mixture Models, Chow-Liu Trees
\end{keywords}

\section{Introduction}

Probabilistic inference has become a core technology in AI,
largely due to developments in graph-theoretic methods for the
representation and manipulation of complex probability
distributions~\citep{pearl:88}.  Whether in their guise as
directed graphs (Bayesian networks) or as undirected graphs (Markov
random fields), \emph{probabilistic graphical models} have a number
of virtues as representations of uncertainty and as inference engines.
Graphical models allow a separation between qualitative, structural
aspects of uncertain knowledge and the quantitative, parametric aspects
of uncertainty...\\

\section{Data}

\subsection{Data set}

\lipsum[1-2]

\begin{table}
\label{tab:1}
\centering
\begin{tabular}{ |c|c|c| }
 \hline
 cell1 & cell2 & cell3 \\
 cell4 & cell5 & cell6 \\
 cell7 & cell8 & cell9 \\
 \hline
\end{tabular}
\caption{My table.}
\end{table}

\lipsum[1]

\subsection{Exploratory data analysis}

\lipsum[1-2]

\section{Methodology}

\lipsum[1-2]

\begin{algorithm}[H]
\begin{algorithmic}[1]
\REQUIRE
$k$: the number of clusters; \\
$D$: a data set containing n objects.
\ENSURE A set of $k$ clusters.

\STATE arbitrarily choose $k$ objects in $D$ as the initial representative objects or seeds;
\REPEAT
 \STATE assign each remaining object to the cluster with the nearest representative object;
 \STATE randomly select a nonrepresentative object, $o_{random}$;
 \STATE compute the total cost $S$ of swapping representative object $o_j$ with $o_random$;
 \IF{$S < 0$}
     \STATE swap $o_j$ with $o_{random}$  to form the new set of $k$ representative objects;
 \ENDIF
\UNTIL no change;
\end{algorithmic}
\caption{k-medoids. PAM, a k-medoids algorithm for partitioning based on medoid
or central objects.}
\label{alg:1}
\end{algorithm}

\section{Results}

\lipsum[1-2]

\section{Conclusions}

\lipsum[1]



% Acknowledgements should go at the end, before appendices and references

\acks{We would like to thank... }

% Manual newpage inserted to improve layout of sample file - not
% needed in general before appendices/bibliography.

\bibliography{bibliography}

\appendix
\section*{Appendix A.}
\label{app:theorem}

% Note: in this sample, the section number is hard-coded in. Following
% proper LaTeX conventions, it should properly be coded as a reference:

%In this appendix we prove the following theorem from
%Section~\ref{sec:textree-generalization}:

In this appendix we prove the following theorem from
Section~6.2:

\noindent
{\bf Theorem} {\it Let $u,v,w$ be discrete variables such that $v, w$ do
not co-occur with $u$ (i.e., $u\neq0\;\Rightarrow \;v=w=0$ in a given
dataset $\dataset$). Let $N_{v0},N_{w0}$ be the number of data points for
which $v=0, w=0$ respectively, and let $I_{uv},I_{uw}$ be the
respective empirical mutual information values based on the sample
$\dataset$. Then
\begin{align*}
N_{v0} \;>\; N_{w0}\;\;\Rightarrow\;\;I_{uv} \;\leq\;I_{uw}
\end{align*}
with equality only if $u$ is identically 0.} \hfill\BlackBox

\noindent
{\bf Proof}. We use the notation:
\begin{align*}
P_v(i) \;=\;\frac{N_v^i}{N},\;\;\;i \neq 0;\;\;\;
P_{v0}\;\equiv\;P_v(0)\; = \;1 - \sum_{i\neq 0}P_v(i).
\end{align*}
These values represent the (empirical) probabilities of $v$
taking value $i\neq 0$ and 0 respectively.  Entropies will be denoted
by $H$. We aim to show that $\fracpartial{I_{uv}}{P_{v0}} < 0$....\\




\end{document}
